% =====================================================================
% chapters/02_Related_work.tex — rendu LaTeX de ../traductions/02_related_work.en_brut.txt (23 septembre 2026)
% Section 2 de la VERSION FRANÇAISE de l'article court. \input depuis main.tex, après 01_Introduction.tex.
% Ce fichier DÉFINIT \label{sec:related}, visé par la carte de lecture de la section 1.4.
% Citations : mcfadden1974conditional, benakiva1985discrete, train2009discrete, verplanken1997habit,
%   garling2003habitual, adam2025survey, park2023generative, liu2024toward, bougie2025citysim,
%   horni2016matsim, lammer2026gta, alves2026evaluating, argyle2023outofone, meister2024benchmarking,
%   baronchelli2025emergent, bintareaf2026silica, wei2022chain, sprague2024cot.
% Rendu mot pour mot de la traduction fournie par l'auteur.
% =====================================================================

\section{Travaux connexes}
\label{sec:related}

Trois courants de recherche se rejoignent dans cet article, et aucun n'évalue si la délibération verbalisée aide les agents à reproduire les parts modales d'un territoire réel.

\subsection{Choix discret et filtres de perception}
\label{sec:discrete-choice}

Depuis cinquante ans, la recherche en transport prédit le choix modal à l'aide de modèles de choix discret, statistiquement robustes et comportementalement rigides. Un tel modèle attribue à chaque mode une probabilité calculée à partir du temps de trajet, du coût et des caractéristiques du voyageur \citep{mcfadden1974conditional,benakiva1985discrete}. Les analystes l'ajustent aux enquêtes de mobilité, dans lesquelles les résidents décrivent chaque trajet d'une journée, de sorte qu'il hérite de leurs parts modales \citep{train2009discrete}. Sa règle de décision, cependant, est une maximisation sur des distributions fixes de préférences.

La recherche comportementale documente des départs qu'aucune enquête n'enregistre. L'habitude rend un choix modal répété automatique. Le voyageur acquiert alors moins d'informations sur les alternatives \citep{verplanken1997habit,garling2003habitual}. \citet{adam2025survey} ont interrogé 650 personnes et ont constaté que les automobilistes réguliers sous-estimaient le prix d'une voiture. Comme un modèle basé sur des enquêtes ne peut intégrer de tels filtres de perception, les agents génératifs sont proposés pour combler cette lacune.

\subsection{Agents génératifs dans la simulation de mobilité}
\label{sec:generative-mobility}

Les modèles de langage ont fait leur apparition dans la simulation de mobilité en tant que composante délibérative absente des simulateurs précédents. \citet{park2023generative} placent le modèle au sein de l'agent, avec une mémoire épisodique et une capacité de réflexion. \citet{liu2024toward} transposent cette approche à la demande de transport, sous forme d'hybride avec des composantes établies. CitySim \citep{bougie2025citysim} déploie ces agents à l'échelle d'une ville et les valide en fonction de l'utilisation du temps. Aucune expérience ne confronte une répartition modale simulée à une répartition modale observée. Dans les simulateurs antérieurs tels que MATSim \citep{horni2016matsim}, l'équilibre résulte de la notation des plans plutôt que de la délibération d'un agent.

GTA \citep{lammer2026gta} est le système publié le plus proche en termes d'évaluation, et trois de ses choix limitent ce qu'il peut afficher. Chaque agent renvoie un mode par trajet plutôt qu'une distribution, donc aucune métrique distributionnelle ne s'applique. Ses jours sont également indépendants. Ses auteurs mentionnent la mémoire multi-jours comme travaux futurs. Nous ajoutons une ligne de base sous les agents et quatre modèles ajustés sur l'enquête au-dessus. Notre métrique évalue une distribution complète, et non un mode par trajet. Nos agents conservent les événements d'un jour à l'autre.

\citet{alves2026evaluating} sont les plus proches du point de vue de la plateforme, bien que leur évaluation ne repose sur aucune vérité humaine. Ils couplent la plateforme GAMA à un module de modélisation du langage doté d'une mémoire persistante, qui détermine, en cas de perturbation, si un agent doit replanifier. Ils évaluent l'adaptabilité par rapport à une base de référence fondée sur des règles, et non par rapport au comportement observé.

\subsection{Alignement distributionnel des populations de modèles de langage}
\label{sec:alignment}

Une autre littérature s'interroge sur la possibilité pour les populations de modèles de langage de se substituer à des échantillons humains. \citet{argyle2023outofone} ont conditionné un modèle sur des attributs sociodémographiques et ont retrouvé les schémas de réponse de la sous-population correspondante. \citet{meister2024benchmarking} ont constaté qu'un modèle verbalisant une distribution s'aligne mieux que le même modèle échantillonné de manière répétée. Notre protocole évalue une distribution verbalisée pour cette raison.

Cette même littérature met en garde contre l'interprétation de ces populations comme étant des populations humaines. \citet{baronchelli2025emergent} constatent des biais collectifs dans les populations de modèles décentralisés, biais qu'aucun agent individuel ne présente. Ils refusent leur statut de substitut humain. L'instrument SILICA \citep{bintareaf2026silica} certifie ces populations à l'aide d'une bibliothèque de perturbations, incluant l'ordre des options. Nous adoptons sa grille de lecture et son contrôle de l'ordre des options.

L'utilité de rédiger le raisonnement avant de répondre fait débat, aucun résultat ne permettant de trancher la question pour une distribution de population donnée. Les gains observés pour le raisonnement écrit proviennent de tâches qui admettent une réponse vérifiable \citep{wei2022chain,sprague2024cot}. Nous effectuons donc nos mesures sur une population d'agents de mobilité, dont la boucle est décrite dans la section suivante.
